\subsection{Endpoint Summaries via RDF Data Shapes}\label{sec:summaries}

The HiBISCuS algorithm~\cite{saleem2014hibiscus} proposes a summary approach, later refined by FedUP~\cite{aimonierdavat:hal-04538238}, to improve federated query performance.
These summaries are \acp{KG} that, for each endpoint, characterize the predicates available for \ac{RDF} terms with respect to the endpoint's \emph{authority}\footnote{The authority as defined by the \href{https://datatracker.ietf.org/doc/html/rfc3986}{Uniform Resource Identifier (URI): Generic Syntax} specification.} in its \acp{IRI}.

We propose a different summarization approach, based on \ac{RDF} data shapes.
At an abstract level, HiBISCuS-type summaries already capture structural regularities of endpoint data, and \ac{RDF} shape languages provide a standard mechanism to represent this structure.
We therefore use \ac{RDF} data shapes to obtain a reusable and more expressive alternative to the custom summaries used by HiBISCuS and FedUP.

Independently of the shape language, we abstract a shape using three definitions,
first a target in \cref{def:target},
then the triple constraints on the entities it selects in \cref{def:triple-constraint},
and finally a shape in \cref{def:shape}.

\begin{definition}[Target]\label{def:target}
A target is a function $t : \mathcal{G} \to 2^{\mathcal{I}\mathcal{B}}$ that selects, from a \ac{KG}, the \acp{IRI} and blank nodes to be described, which we call \emph{entities}.
\end{definition}

A target $t$ \emph{can select} the entities $\mathit{sel}(t) = \bigcup_{G \in \mathcal{G}} t(G)$.
This set is known without the data when $t$ selects a set of \acp{IRI}, listed explicitly or expressed by a regular expression,
while it is left to future work for other targets, such as the instances of a class.

\begin{definition}[Triple constraint]\label{def:triple-constraint}
Let $t$ be a target.
A triple constraint for $t$ is a triple $(p, C_O, r)$ constraining, in a \ac{KG} $G$, the triples $(x, p, o) \in G$ whose subject $x$ is in $t(G)$, where $p$ is a predicate,
$C_O$ is a set of options for the object of $p$, each being either
(i) a constraint on the value of the object, such as a single allowed value, or
(ii) a reference $@s'$ to a shape $s'$, requiring the object of $p$ to conform to $s'$,
and $r \in \{\mathit{true}, \mathit{false}\}$ states whether $p$ is required.
\end{definition}

\begin{definition}[Shape]\label{def:shape}
A shape is a triple $s = (t, C_E, \mathit{TC})$, where
$t$ is a target as defined in \cref{def:target},
$C_E$ is a constraint on the entities themselves, such as a regular expression on their \acp{IRI}, or $\top$ when there is none, and
$\mathit{TC} = \{A_1, \dots, A_k\}$ is a set of alternatives, each alternative $A_j$ being a set of triple constraints for $t$ as defined in \cref{def:triple-constraint}.
Given a \ac{KG} $G$, a term $o$ satisfies $C_O$ if it satisfies one of the value constraints of $C_O$, or conforms to $s'$ for one of its references $@s'$.\footnote{Shapes can refer to each other cyclically, for instance a shape for persons whose \texttt{foaf:knows} objects are persons.}
An entity $x \in t(G)$ satisfies a triple constraint $(p, C_O, r)$ if the object of every triple $(x, p, o) \in G$ satisfies $C_O$
and, when $r = \mathit{true}$, $G$ contains at least one such triple.
The entity $x$ conforms to $s$ if it satisfies $C_E$ and, for some alternative $A_j \in \mathit{TC}$, it satisfies every triple constraint of $A_j$
and the object $o$ of every triple $(x, p, o) \in G$ satisfies $C_O$ for some $(p, C_O, r) \in A_j$.
\end{definition}

The \acp{IRI} of the entities conforming to a shape $s = (t, C_E, \mathit{TC})$ are thus in the subset
$\mathit{ent}(s) = \{ x \in \mathit{sel}(t) \cap \mathcal{I} \mid x \text{ satisfies } C_E \} \subseteq \mathit{sel}(t)$
of the entities that $t$ can select, where $\mathit{sel}(t)$ is taken as $\mathcal{I}$ when it is unknown.

To describe the content of each federation member, we use a summary method called \emph{shape index}~\cite{tam2025},
``a mapping between \ac{RDF} document sets and \ac{RDF} data shapes that describes a subweb controlled by a data provider''~\cite{tam2025},
where a subweb is ``a (sub)\ac{DKG} defined by the \ac{KG} derived from a set of \acp{IRI} controlled by a data provider''~\cite{tam2025}.
A shape index is defined as

\begin{equation}\label{eq:shape-index}
  SI = \{ s_1 \mapsto \mathit{IRI}_1, s_2 \mapsto \mathit{IRI}_2, \dots, s_n \mapsto \mathit{IRI}_n \}
\end{equation}

where $s_i$ is a shape and $\mathit{IRI}_i = SI(s_i)$ is a set of \acp{IRI}, and it is \emph{complete} when every shape $s_i \in \mathrm{dom}(SI)$ is closed~\cite{tam2025}.

In a federation, a subweb is a set of federation members, identified by the \acp{IRI} of their SPARQL endpoints.
Each $\mathit{IRI}_i$ is thus the set of all endpoint \acp{IRI} whose \ac{KG} contains triples of an entity conforming to $s_i$.
The mapping $s_i \mapsto \mathit{IRI}_i$ therefore states that the entities selected by the target of $s_i$, together with their triples constrained by the alternatives of $s_i$, are stored in the endpoints of $\mathit{IRI}_i$,
whereas the objects of these triples that must conform to a shape $s_k$ through a triple constraint of kind~(ii) \emph{can be described} by the endpoints of $\mathit{IRI}_k$.
In most cases, $\mathit{IRI}_i$ contains a single endpoint.
It can however contain several, for instance when data following the same schema are partitioned over multiple endpoints and the constraints of $s_i$ do not distinguish the partitions.

We further require the shape index to describe every triple of the endpoints it maps.
Let $\mathit{KG}(e)$ be the \ac{KG} of the endpoint identified by the \ac{IRI} $e$, and
$G_{SI} = \bigcup_{e \in E_{SI}} \mathit{KG}(e)$, with $E_{SI} = \bigcup_{s_i \in \mathrm{dom}(SI)} \mathit{IRI}_i$,
the \ac{KG} of the subweb described by $SI$.
Conformance to the shapes of $SI$ is evaluated over $G_{SI}$,
so an entity whose triples are stored in several endpoints conforms to a shape only if all its triples in the federation satisfy it.
A triple $(x, p, o)$ is \emph{described} by a shape $s_i = (t_i, C_{E,i}, \mathit{TC}_i)$ if $x \in t_i(G_{SI})$, meaning that its subject is selected by $t_i$ over this subweb,
and $s_i$ accepts the predicate $p$, meaning that $(p, C_O, r) \in A$ for some alternative $A \in \mathit{TC}_i$, some $C_O$, and some $r$.
The shape index must satisfy
\begin{equation}\label{eq:coverage}
  \forall e \in E_{SI},\ \forall (x, p, o) \in \mathit{KG}(e),\ \exists s_i \in \mathrm{dom}(SI) :
  e \in \mathit{IRI}_i \land (x, p, o) \text{ is described by } s_i.
\end{equation}
Shape languages also allow open shapes, which accept any predicate not forbidden by a negative constraint.
Open shapes satisfy \cref{eq:coverage} trivially but describe little about the data.
We therefore only consider closed shapes, as captured by \cref{def:shape}, and thus complete shape indexes.
We further assume that the data conforms to the shapes of $SI$, meaning that every entity selected by $t_i$ over $G_{SI}$ conforms to $s_i$, as when a data provider validates its subweb against its shape index.

Closed shapes also determine where the triples of an entity are stored.
Each data provider describes its own subweb,
so an entity whose triples are split over the subwebs of several providers would require a closed shape over data that no single provider controls.
We therefore assume \emph{entity locality}, meaning that the triples of an entity accepted by a shape are stored in the endpoints of that shape,
\begin{equation}\label{eq:locality}
  \forall s_i \in \mathrm{dom}(SI),\ \forall x \in t_i(G_{SI}) :
  \{ (x, p, o) \in G_{SI} \mid s_i \text{ accepts } p \} \subseteq \bigcup_{e \in \mathit{IRI}_i} \mathit{KG}(e).
\end{equation}
The triples of an entity can thus be split over several endpoints of $\mathit{IRI}_i$,
whereas the triples whose subject is an object of this entity conforming to a nested shape can be stored in the endpoints of other shapes.

Shape languages can express exclusive unions between sets of triple constraints, which \cref{def:shape} captures with the alternatives of $\mathit{TC}$.
By \cref{def:shape}, an entity conforming to a shape $s = (t, C_E, \mathit{TC})$ satisfies an alternative $A_j \in \mathit{TC}$,
and each of its triples uses a predicate of a triple constraint of $A_j$.
A shape index therefore describes not only which predicates occur in an endpoint, but also which combinations of predicates a single entity can have.

\lstinputlisting[caption={ShEx example with an exclusive choice of user name representation.},label={lst:user-alt}]{code/user_alt.shex}

For example, consider the ShEx shape in \cref{lst:user-alt}.
It describes users that must have a \texttt{foaf:mbox} and either:
\begin{itemize}
  \item a \texttt{foaf:name}, or
  \item a \texttt{foaf:givenName} and a \texttt{foaf:familyName}.
\end{itemize}

Since \texttt{ex:UserShape} is closed, a user conforming to it follows exactly one of these alternatives.
Hence, if \texttt{ex:UserShape} is mapped to a set of endpoint \acp{IRI} in a shape index, the index describes that no user stored in these endpoints has both a \texttt{foaf:name} and a \texttt{foaf:givenName}.
Different users may nevertheless follow different alternatives, so the endpoints can contain users of both kinds.

Beyond the combinations of predicates that an entity can have, triple constraints also describe the objects of these predicates,
either by constraining their value, kind~(i), or by requiring them to conform to another shape, kind~(ii).



\lstinputlisting[caption={ShEx example using object constraint.},label={lst:shex-object-constraint}]{code/object_constraint.shex}

For instance, in \cref{lst:shex-object-constraint}, the triple constraints
\texttt{(rdf:type, [foaf:Person foaf:User], true)}, \texttt{(foaf:name, LITERAL, true)}, and
\texttt{(foaf:mbox, IRI, false)} of \texttt{ex:UserShape}, and
\texttt{(foaf:name, xsd:string, true)} of \texttt{ex:Hobby}, are of kind~(i),
as they constrain the value of the object,
meaning that users are of type \texttt{foaf:Person} or \texttt{foaf:User}, their name is a literal, their mailbox is an \ac{IRI}, and hobbies have a name of type \texttt{xsd:string}.
In contrast, \texttt{(ex:hobby, @ex:Hobby, true)} is of kind~(ii),
as it requires the object of \texttt{ex:hobby} to conform to \texttt{ex:Hobby}.
All these triple constraints are required, except \texttt{(foaf:mbox, IRI, false)},
since the \texttt{?} of \cref{lst:shex-object-constraint} makes the mailbox optional,
so a user may have no mailbox, but any mailbox it has is an \ac{IRI}.
Mapped to a set of endpoint \acp{IRI} in a shape index, these triple constraints thus describe not only the predicates of the users stored in these endpoints, but also their objects.
